% Generated draft: requires mathematical and source review.
\documentclass[11pt]{article}
\usepackage[margin=25mm]{geometry}
\usepackage{fontspec}
\setmainfont{Latin Modern Roman}
\usepackage{amsmath,amssymb,mathtools,unicode-math}
\setmathfont{Latin Modern Math}
\usepackage{enumitem,needspace,xurl,hyperref}
\setlength{\parindent}{0pt}
\setlength{\parskip}{5pt}
\setlength{\emergencystretch}{3em}
\newcommand{\N}{\mathbb{N}}
\newcommand{\Z}{\mathbb{Z}}
\newcommand{\Q}{\mathbb{Q}}
\newcommand{\R}{\mathbb{R}}
\newcommand{\C}{\mathbb{C}}
\newcommand{\F}{\mathbb{F}}
\newcommand{\A}{\mathbb{A}}
\newcommand{\PP}{\mathbb{P}}
\newcommand{\E}{\mathbb{E}}
\newcommand{\eps}{\varepsilon}
\newcommand{\dd}{\,\mathrm{d}}
\begin{document}
\section*{UnsolvedMath ID: 800608}
\textbf{Source problem code:} catalogue entry
\subsection{Definitions and mathematical statement}
For a topological space \(X\), let \(C_p(X)\) denote the space of continuous functions \(X\to\R\), equipped with the topology of pointwise convergence, and let \(\mathbf{0}\) denote its zero function. Using the usual convergence-set convention, let \(\Gamma_{\mathbf{0}}\) consist of the infinite sets \(A\subseteq C_p(X)\setminus\{\mathbf{0}\}\) such that every neighborhood of \(\mathbf{0}\) contains all but finitely many members of \(A\). Equivalently, for each finite \(F\subseteq X\) and each \(\varepsilon>0\), all but finitely many \(f\in A\) satisfy \(|f(x)|<\varepsilon\) for every \(x\in F\).

The property \(\mathsf{S}_1(\Gamma_{\mathbf{0}},\Gamma_{\mathbf{0}})\) means that for every sequence \((A_n)_{n\in\N}\) of members of \(\Gamma_{\mathbf{0}}\), there are choices \(f_n\in A_n\) such that \(\{f_n:n\in\N\}\in\Gamma_{\mathbf{0}}\). These definitions explain the notation under a standard convention; the supplied record does not explicitly give the definitions or the underlying category of spaces.

For \(f,g\in\N^{\N}\), write \(f\leq^*g\) if \(f(n)\leq g(n)\) for all but finitely many \(n\). The bounding number is
\[
\mathfrak b=\min\bigl\{|B|:B\subseteq\N^{\N},\ (\forall g\in\N^{\N})(\exists f\in B)\ f\not\leq^*g\bigr\}.
\]

The supplied Problem 36 asks: assuming
\[
\kappa<\mathfrak b
\quad\text{and}\quad
(\forall\alpha<\lambda)\;
C_p(X_\alpha)\text{ satisfies }\mathsf{S}_1(\Gamma_{\mathbf{0}},\Gamma_{\mathbf{0}}),
\]
must
\[
C_p\!\left(\bigcup_{\alpha<\kappa}X_\alpha\right)
\text{ satisfy }\mathsf{S}_1(\Gamma_{\mathbf{0}},\Gamma_{\mathbf{0}})?
\]
The hypothesis is indexed by \(\lambda\), whereas the union is indexed by \(\kappa\). The supplied record neither defines \(\lambda\) nor relates it to \(\kappa\); this discrepancy is retained and prevents a unique precise interpretation. It also does not specify the class of spaces \(X_\alpha\) or the topology on their union. If they are subspaces of a common topological space, a subspace topology on the union would supply such a topology, but this ambient-space hypothesis is not supplied. No identification of \(\lambda\) with \(\kappa\), or choice of ambient category, is asserted here.
\subsection{Short English statement}
The question asks whether taking a union of fewer than \(\mathfrak b\) spaces preserves the stated selection property of their pointwise continuous-function spaces. The source uses different index bounds in the hypothesis and conclusion and leaves the topology of the union unspecified, so this intended preservation question requires clarification. The registry labels it open; that label is supplied information, not an independently verified assessment.
\subsection{Sources}
\begin{itemize}
\item Local UnsolvedMath registry, permanent ID \(800608\), exact source problem code \texttt{catalogue entry}; upstream ID \(800608\), revision \texttt{d371ee917ac9647fbb06cd6fc529deec6f3dda82}. The supplied record attributes the problem to Tsaban, gives the year 2006, and labels its status as open.
\item Tsaban, \emph{Selection Principles and Special Sets of Reals: Open Problems} (2006), Problem 36. Supplied source URL: \url{https://arxiv.org/abs/math/0609601}. The record reports access on 2026-10-16 and extraction by \texttt{arxiv\_latex}; these are supplied provenance details, not a claim of independent source access or verification.
\item The supplied title and background attach the attribution/context label \texttt{[FremlinSeq]}. No bibliographic expansion of this label is supplied, so its precise reference and attribution remain unresolved.
\end{itemize}
\end{document}
